\section{Constant-depth threshold circuits and a lower-bound barrier}
\label{sec:threshold-circuits}

The coefficient-independent halfspace construction also clarifies which
explicit targets could disprove the general unrestricted sample
conjecture. Here a threshold circuit may have arbitrary real weights,
and its size counts wires, including input and constant wires.

\begin{lemma}[Threshold circuits as bounded-weight tanh networks]
\label{lem:threshold-circuit-tanh}
There is a universal constant $a$ such that every depth-$h$, $h\ge1$, threshold
circuit of wire size at most $s$ on $d$ Boolean inputs has a tanh
realization with linear readout, unit global margin, depth at most
$2h$, and both width and maximum absolute physical weight bounded by
$(2+d+s)^a$. It uses no biases after the augmented input. The bound is
independent of the original real coefficients and of every original
gate's geometric margin.
\end{lemma}

\begin{proof}
Replace each threshold gate by its depth-two polynomial-size,
polynomial-integer-weight simulation from
Lemma~\ref{lem:goldmann-karpinski-simulation}. Each gate has at most
$s$ inputs. The total wire count, coefficient bound, and maximum
weighted fan-in of the resulting circuit are bounded by a universal
polynomial in $d+s$. Its depth is at most $2h$. Subdivide skip wires by
identity threshold gates to make the circuit layered; since $h\le s$,
this costs only another polynomial factor. Pad with identity gates as
needed so that the output is at a common depth $q=2h\ge2$.

Every gate can be taken to have integer affine form with absolute value
at least one on Boolean inputs. Let $B\ge1$ bound the sum of the
absolute coefficients of any such affine form, including its constant.
Put $\delta=1/(4B)$ and $\lambda=2\log(8B)$. Replace every gate except
the final output gate by tanh applied to $\lambda$ times its affine
form. If all preceding features approximate their intended signs to
accuracy $\delta$, the affine perturbation is at most $B\delta=1/4$.
Its signed value is therefore at least $3/4$, and the new feature error
is at most
\[
 2e^{-3\lambda/2}=\frac{2}{(8B)^3}\le\delta.
\]
This inductive accuracy bound does not deteriorate with depth.

To implement constants without later biases, add a constant feature
equal to $1/2$ at every hidden layer. The first such feature has bias
$\operatorname{atanh}(1/2)$; each later one receives weight
$2\operatorname{atanh}(1/2)$ from the preceding constant feature.
An affine constant $b$ is then supplied by weight $2\lambda b$ to
this feature. Thus all constants are exact. The last threshold is
replaced by twice its affine form, as a linear readout; its signed
value is at least $3/2$. Uniform zero-padding of the hidden layers
preserves the function. Every used parameter is polynomial in $d+s$,
as asserted.
\end{proof}

\begin{corollary}[Precision-independent threshold-circuit learning]
\label{cor:threshold-circuit-learning}
Fix $\zeta>0$. A target with a depth-$h$, wire-size-$s$ threshold circuit
has unrestricted Gaussian posterior sample complexity
\[
 (2+d+s)^{O_\zeta(h)}
 \bigl(L+\varepsilon^{-1}+\log(2/\delta)\bigr)^{O_\zeta(1)}
\]
at every learner depth $L\ge2h$, simultaneously for every
$N\ge C_L(f)$ and every $0<\sigma^2\le N^{-(1+\zeta)}$.
No original coefficient magnitude or geometric margin enters the bound.
In particular, for each fixed $h$, this is polynomial in the circuit
size, input dimension, learner depth, and statistical parameters.
\end{corollary}

\begin{proof}
Apply the fixed-intrinsic-depth theorem and its quantitative corollary
below to Lemma~\ref{lem:threshold-circuit-tanh}, with intrinsic depth
$D_0=2h$ and polynomial bounds on the tanh width and weights.
\end{proof}

Recall that nonuniform $\mathsf{TC}^0$ consists of Boolean families
having polynomial-size threshold circuits of some fixed depth.
Consequently every such family enjoys a polynomial sample guarantee
under the present posterior at all sufficiently large learner depths.
This gives the following precise limitation on a negative-result route.

\begin{corollary}[A circuit-complexity consequence of a counterexample]
\label{cor:tc0-counterexample-barrier}
Fix $\zeta>0$. Suppose a family $(f_d)$, at learner depths $L(d)\to\infty$,
requires more than every polynomial in $d,L(d)$ examples for some fixed
positive error and confidence targets, under the present unrestricted
posterior and the all-admissible-width guarantee with
$\gamma\ge1+\zeta$. Then $(f_d)\notin\mathsf{TC}^0$.
In particular, a family in $\mathsf{NC}^1$ with this property would
separate $\mathsf{NC}^1$ from $\mathsf{TC}^0$.
\end{corollary}

\begin{proof}
If the family were in $\mathsf{TC}^0$, it would have fixed depth $h$
and wire size polynomial in $d$. Eventually $L(d)\ge2h$, so
Corollary~\ref{cor:threshold-circuit-learning} would supply a polynomial
sample bound simultaneously for every admissible width and stated
Gaussian scale, a contradiction.
\end{proof}

This is a logical consequence, not a claimed circuit separation or a
proof that unrestricted learning is always polynomial. It also does not
cover exponents approaching one with dimension, fixed learner depths
below the alternative representation's depth, or the mean-field case.
